% Generated draft: requires mathematical and source review.
\documentclass[11pt]{article}
\usepackage[margin=25mm]{geometry}
\usepackage{fontspec}
\setmainfont{Latin Modern Roman}
\usepackage{amsmath,amssymb,mathtools,unicode-math}
\setmathfont{Latin Modern Math}
\usepackage{enumitem,needspace,xurl,hyperref}
\setlength{\parindent}{0pt}
\setlength{\parskip}{5pt}
\setlength{\emergencystretch}{3em}
\newcommand{\N}{\mathbb{N}}
\newcommand{\Z}{\mathbb{Z}}
\newcommand{\Q}{\mathbb{Q}}
\newcommand{\R}{\mathbb{R}}
\newcommand{\C}{\mathbb{C}}
\newcommand{\F}{\mathbb{F}}
\newcommand{\A}{\mathbb{A}}
\newcommand{\PP}{\mathbb{P}}
\newcommand{\E}{\mathbb{E}}
\newcommand{\eps}{\varepsilon}
\newcommand{\dd}{\,\mathrm{d}}
\begin{document}
\section*{UnsolvedMath ID: 7200005}
\textbf{Source problem code:} catalogue entry
\subsection{Definitions and mathematical statement}
In a standard algebraic formulation, let $k$ be a commutative ring and $A$ a unital associative $k$-algebra. Its Hochschild cochains are $C^q(A,A)=\operatorname{Hom}_k(A^{\otimes q},A)$ for $q\geq0$, with differential
\[
\begin{aligned}
(bf)(a_1,\ldots,a_{q+1})={}&a_1f(a_2,\ldots,a_{q+1})\\
&+\sum_{i=1}^{q}(-1)^i f(a_1,\ldots,a_i a_{i+1},\ldots,a_{q+1})\\
&+(-1)^{q+1}f(a_1,\ldots,a_q)a_{q+1}.
\end{aligned}
\]
Write $HH^*(A,A)=H^*(C^*(A,A),b)$.

Deligne's conjecture asks for an $E_2$-algebra structure on this cochain complex: construct an action of a differential graded operad modeling chains on the little two-disks operad, with the usual conversion from chain degrees to cohomological degrees. Little two-disks operations are configurations of disjoint disks embedded by translations and positive dilations in one disk, composed by insertion. The induced operations on $HH^*(A,A)$ should recover the Hochschild cup product and Gerstenhaber bracket of degree $-1$, with their Gerstenhaber-algebra identities. The target is this chain-level lift, not merely the existence of those cohomology operations.

The supplied record names the conjecture but does not give a ground-ring, normalization, or differential-graded-algebra convention; the displayed ordinary-algebra setup makes one standard version explicit. Its summary reports multiple proofs. Its additional claim that the $E_2$ structure is unique up to homotopy requires specified models and compatibility/naturality data and is not appended as a consequence of the existence conjecture.
\subsection{Short English statement}
Give Hochschild cochains an action of a chain model of the little two-disks operad that lifts the cup product and Gerstenhaber bracket on Hochschild cohomology. The supplied record reports this existence conjecture solved; blanket uniqueness is a separate assertion.
\subsection{Sources}
\begin{itemize}
\item Wikipedia contributors, \emph{List of unsolved problems in mathematics}, Geometry bullet 5, pinned revision 1366636767, \url{https://en.wikipedia.org/w/index.php?title=List_of_unsolved_problems_in_mathematics&oldid=1366636767#Geometry}; accessed 2026-07-29, MediaWiki source extraction. Wikipedia text is CC BY-SA 4.0, with attribution/share-alike handling required for release.
\item Supplied proof attributions: McClure--Smith (2002--2006), \emph{Forum Mathematicum}/\emph{Contemporary Mathematics}, including the title ``A solution of Deligne's conjecture''; Tamarkin (1998/2015), \emph{Advances in Mathematics}; Kontsevich--Soibelman, ``Deformation quantization and the Koszul operad'' (2000), as titled in the supplied summary; Voronov (2000), Swiss-cheese-type operads; Berger--Fresse (2004), surjection operad. The record supplies no precise arXiv identifiers or complete publication data for these entries.
\item Local registry \texttt{data/unsolvedmath/problems.json}, permanent UnsolvedMath ID $7200005$, source code catalogue entry, \emph{Deligne's conjecture on Hochschild cohomology about the operadic structure on Hochschild cochain complex}. Dataset attribution: Ulam AI and the UnsolvedMath Contributors.
\end{itemize}
\end{document}
